\chapter{Introducción al desarrollo realizado}
\label{chap:desenrolo}

\section{Introducción}


\lettrine{E}{n} este capítulo se presenta una descripción del desarrollo realizado, los objetivos perseguidos y la metodología aplicada durante la implementación. El propósito es ofrecer un vistazo general de los artefactos construidos y de las decisiones más relevantes tomadas durante el desarrollo del proyecto.

\section{Arquitectura de la Solución}


\vspace{1em}

{ \color{red} PENDIENTE una vez avanzado más el proyecto }

\section{Metodología e iteraciones}

El proceso de desarrollo se organizó mediante iteraciones cortas con objetivos funcionales concretos, siguiendo las ideas de Scrum adaptadas al contexto de trabajo individual. Aunque no se aplicó una implementación completa y formal de Scrum (debido al tamaño reducido del equipo), se adoptaron sus pilares fundamentales: planificación por sprints, revisiónes y \emph{feedback} contínuo.

En cuanto a roles, el autor del \acrshort{tfm} asumió las responsabilidades de desarrollo y encargado de definición funcional, mientras que el tutor actuó como responsable de supervisión y de aceptación.

Se establecieron reuniones de revisión mensuales con el tutor en las que se presentaron los artefactos construidos (prototipos, endpoints, pantallas, pruebas y documentación). Estas sesiones sirvieron para validar funcionalidades, priorizar tareas y ajustar el dominio del \acrshort{tfm} y el alcance de las siguientes iteraciones.

El proyecto se estructuró en una serie de sprints. De una forma resumida, las fases necesarias para realizar este \acrshort{tfm} fueron: una fase inicial de contextualización y definición del valor del proyecto; la creación del esqueleto y la arquitectura base de la aplicación móvil para facilitar el desarrollo posterior; la normalización y actualización del modelo de recursos y organizaciones manteniendo compatibilidad con las entidades Equipo y Vehículo; la implementación de la gestión de emergencias y asignaciones básicas; la introducción de un esquema de logs operativos y la migración o compatibilidad con tablas legacy; la integración de notificaciones móviles mediante \acrshort{fcm} y el flujo de aceptación de asignaciones; y, finalmente, la implementación de capacidades semiautomáticas junto a un motor de reglas sencillo para generar propuestas automáticamente.

Para la planificación y el detalle de cada sprint se describen con mayor extensión en el capítulo de planificación (capítulo \ref{sec:sprints}).
